
\lussec 8. $\mib\ZA$ 与 ${\mib c}_{\bf fl}$ 的计算策略

只要恰当地选取相关的重正化常数与改善系数，
PCAC 关系 (6.32) 就成立。
如下文所解释的，
反过来，通过要求该关系在运动学参数空间的两个点上精确成立，
可以在分别为 $a^2$ 阶与 $a$ 阶精度内确定常数 $\ZA$
与系数 $\cfl$。

按此思路通过数值模拟计算 $\ZA$ 与 $\cfl$，
假定改善系数 $\csw,\ca$ 已经算出。
空间方向应施加周期性边界条件，
除此之外不再作其他假设。
特别地，无需指明时间方向的边界条件，
因为算得的 $\ZA$ 与 $\cfl$ 的值
只在格距效应的水平上依赖于它们。


\lussubsec 8.1 PCAC 关系的积分形式

不幸的是，当把重正化的 PCAC 关系
用未改进裸场的关联函数展开时，
O($a$) 改进会导致项数膨胀。
不过，将等式对围绕 $y_0$ 的时段 $x_0$
以及空间格点的所有 $\vec{x}$ 求和，
可以实现一点简化。
进入该关系的裸关联函数之一便是
\equation{
  \CPP{rs}(t,d)=\sum_{x_0=y_0-d}^{y_0+d}\sum_{\vec{x}}
  \langle\Pax^{rs}(x)\Pax_t^{sr}(y)\rangle,
  \enum
}
另外还有 $6$ 个这样的关联函数，
其中 $P^{rs}$ 被换成
\equation{
  \tilde{P}^{rs},\hat{P}^{rs},Q^{rs},\ring{\partial}_{\mu}A_{\mu}^{rs},
  \ring{\partial}_{\mu}\tilde{A}_{\mu}^{rs}
  \kern0.5em\hbox{or}\kern0.5em
  \drvstar{\mu}\drv{\mu}P^{rs}.
  \enum
}
对此清单中的最后三个场，
对 $x_0$ 的求和按例如
\equation{
  \CAP{rs}(t,d)=
  \frac{1}{2}\sum_{\vec{x}}\bigl\{
  \langle(\left.A^{rs}_0(x)\right|_{x_0=y_0+d+1}+
          \left.A^{rs}_0(x)\right|_{x_0=y_0+d})
  P^{sr}_t(y)\rangle
  \noenum
  \nexteq{1.0ex}
  {\phantom{\CAP{rs}(t,d)=
  \frac{1}{2}\sum_{\vec{x}}}}\kern-0.65em
  -\langle(\left.A^{rs}_0(x)\right|_{x_0=y_0-d-1}+
           \left.A^{rs}_0(x)\right|_{x_0=y_0-d})
  P^{sr}_t(y)\rangle
  \bigr\},
  \enum
}
的方式化归为求和范围边界处各项之和。

为了记号方便，现在引入缩写
\equation{
  \ZAh{rs}=\ZA\left(1+b_A\mq{rs}+\bar{b}_A\tr\Mq\right),
  \enum
  \nexteq{2.5ex}
  \ZPh{rs}=\ZP\left(1+b_P\mq{rs}+\bar{b}_P\tr\Mq\right).
  \enum
}
通过
\equation{
  m_{rs}={\ZPh{rs}\over\ZAh{rs}}(\mR{r}+\mR{s}),
  \enum
}
定义裸流夸克质量之和后，
积分形式的 PCAC 关系变为
\equation{
  \ZAh{rs}\bigl\{\CAP{rs}+
  \tilde{c}_A\CAtP{rs}+c_A\CdPP{rs}
  -m_{rs}\bigl(\CPP{rs}+\tilde{c}_P\CPtP{rs}\bigr)\bigr\}
  \noenum
  \nexteq{2.0ex}
  \qquad+\CPtP{rs}+\hat{c}_A\CAtP{rs}-2\cfl\CPhP{rs}
  -b_{\chi}\frac{1}{2}(\mq{r}-\mq{s})\CQP{rs}=\rmO(a^2).
  \enum
}
这个等式看起来仍然相当复杂，
但应当记住：
它是若干可计算的关联函数之间的线性关系，
必须对所有流时间 $t>0$、所有时间段 $d\geq0$
以及所有夸克味道 $r\neq s$ 成立。
因此未知系数受到该恒等式的强烈限制。

当 $d$ 大于梯度流的光滑范围 $\sqrt{8t}$ 时，
$\CAtP{rs}$ 按高斯型衰减，
相对式 (8.7) 中其余各项迅速变得可忽略。
而且此时花括号中各项之和以及等式中所有其他项
实际上都与 $d$ 无关。
除非另有说明，
下文只考虑这一 $d$ 范围，
并略去正比于 $\CAtP{rs}$ 的两项。


\lussubsec 8.2 上-下味道道

在上、下夸克质量简并的 QCD 中，
上下道的 PCAC 关系取如下形式：
\equation{
  \ZAh{ud}\bigl\{\CAP{ud}+c_A\CdPP{ud}-m_{ud}\CPP{ud}\bigr\}
  -2\cfl\CPhP{ud}
  =-\bigl(1-\ZAh{ud}\tilde{c}_Pm_{ud}\bigr)\CPtP{ud}
  \enum
}
（为简单起见，从此处起省略 O($a^2$) 误差项）。
裸质量组合 $m_{ud}$ 可以照常由流时间为零处的
PCAC 关系确定，并可在此视为已知。
在两个流时间 $t$ 值处计算式 (8.8) 中的关联函数，
便得到关于比值
\equation{
  {\ZAh{ud}\over1-\ZAh{ud}\tilde{c}_Pm_{ud}}
  \quad\hbox{and}\quad
  {\cfl\over1-\ZAh{ud}\tilde{c}_Pm_{ud}}.
  \enum
}
的两个线性方程。
显然，只要所选的流时间处于标度区
并以物理单位固定，
算得的比值就不依赖于它们，
误差分别不超过 $a^2$ 阶与 $a$ 阶。

当夸克质量接近物理值且格距 $a\leq0.1$ fm 时，
因子 $1-\ZAh{ud}\tilde{c}_Pm_{ud}$
与 1 相差不到百分之一。
该修正项的大小可以通过代入已知的 $m_{ud}$
及 $\tilde{c}_{P}$ 的树图值 (6.22) 来估计。
于是 $\ZAh{ud}$ 与 $\cfl$ 可以在这一小修正因子的精度内算出。
注意：通过改进轴矢流的真空到 $\pi$ 介子矩阵元确定的
裸 $\pi$ 介子衰变常数，
应当用依赖质量的因子 $\ZAh{ud}$ 而非 $\ZA$ 重正化。
不过差别很小，
可以用树图值 $b_A=1$ 作出估计。


\lussubsec 8.3 上-奇异味道道

此道中 PCAC 关系
\equation{
  \ZAh{us}\bigl\{\CAP{us}+c_A\CdPP{us}-m_{us}\CPP{us}\bigr\}
  -2\cfl\CPhP{us}-b_{\chi}\frac{1}{2}(\mq{u}-\mq{s})\CQP{us}
  \noenum
  \nexteq{2.0ex}
  \qquad
  =-\bigl(1-\ZAh{us}\tilde{c}_Pm_{us}\bigr)\CPtP{us}
  \enum
}
出现的附加项正比于质量差 $\mq{u}-\mq{s}$ 的平方，
因此实际中很可能可以忽略。
为了能够确定比值 (8.9)（把 $ud$ 换成 $us$），
除此之外还须在三个流时间 $t$ 值处计算式 (8.10)。

对于 $a\leq0.1$ fm 及物理夸克质量，
修正因子 $1-\ZAh{us}\tilde{c}_Pm_{us}$
与 1 的偏离可达百分之一左右。
与上下道所得结果的比较，
即
\equation{
  {\ZAh{us}(1-\ZAh{ud}\tilde{c}_Pm_{ud})
  \over\ZAh{ud}(1-\ZAh{us}\tilde{c}_Pm_{us})}
  =1+b_A\frac{1}{2}(\mq{s}-\mq{d})+\ZA\tilde{c}_P(m_{us}-m_{ud}),
  \enum
}
不幸只能给出系数 $b_A$ 与 $\tilde{c}_P$
的某种线性组合的估计。
